\documentclass{article}

\usepackage{graphicx} % Paquete para incluir imágenes
\usepackage[utf8]{inputenc} % Para caracteres especiales
\usepackage[spanish]{babel} % Para idioma español
\usepackage{xcolor} % Para colores
\usepackage{amsmath} % Para ecuaciones
\usepackage{amssymb} % Para símbolos matemáticos adicionales (incluye \therefore)
\usepackage{fancyhdr} % Para encabezados y pies de página
%\usepackage{geometry} % Para ajustar márgenes
\usepackage{caption} % Para usar \captionof fuera de float
\usepackage{float}
\usepackage{multirow}
\usepackage{enumitem}
\usepackage{lastpage}% Para el número de la última página
\usepackage{hyperref}
\hypersetup{
    colorlinks=true,
    linkcolor=black,
    citecolor=black,
    urlcolor=black,
    pdfborder={0 0 0}
}

% Configuración de márgenes
\usepackage[a4paper,
  top=1.9cm,
  bottom=3.67cm,
  left=1.9cm,
  right=1.32cm
]{geometry}


% Configuración de encabezados y pies de página
\pagestyle{fancy}
\fancyhf{} % Limpia encabezados y pies de página

% Encabezado
\fancyhead[L]{\scriptsize UNCuyo - Ing. Mecatrónica\\Mendoza - Argentina}
\fancyhead[C]{\scriptsize\textbf{Espacio Curricular:  \textbf{AUTÓMATAS Y CONTROL DISCRETO (Año 2024)}\\
\textbf{PROYECTO GLOBAL INTEGRADOR}}}
\fancyhead[R]{\scriptsize Peña Lautaro, Peralta Bruno\\19/07/2025}

% Pie de página
\fancyfoot[C]{Pág. \thepage\ de \pageref{LastPage}}

\setlength{\headheight}{16.9528pt}  % importante si usás fancyhdr con varios renglones
\setlength{\headsep}{0.5cm}    % separación entre cabecera y texto

% \fancyhead[L]{Automática y Máquinas Eléctricas} % Encabezado izquierdo
% \fancyhead[C]{2025} % Encabezado central
% \fancyhead[R]{Peña Lautaro, Peralta Bruno} % Encabezado derecho
% \fancyfoot[C]{\thepage} % Número de página centrado en el pie de página

\begin{document}

% Portada
\begin{titlepage}
    \centering

    % % Logo
    % \includegraphics[width=0.8\textwidth]{imagenes/fing.jpeg}\par\vspace{1cm}
    % \vspace{2cm}
    % % Título
    % {\Huge \textbf{Proyecto Global Integrador:}\\[0.5cm]
    % \textbf{Control de Accionamiento de CA con}\\
    % \textbf{Motor Sincrónico de Imanes Permanentes}\par}\vspace{2cm}
    
    % Autores
    {\Large Peña Lautaro - 13099\\
    Peralta Bruno - 13220\par}\vspace{0.5cm}
    
    % Profesor titular
    {Año 2024\par}
\end{titlepage}

% Índice
\newpage
\tableofcontents
\newpage

% Contenido del documento
\section{Resumen}
Este documento presenta el diseño y análisis de un sistema 
de control para un accionamiento eléctrico 
de corriente alterna (CA) compuesto por una máquina eléctrica sincrónica con excitación
por imanes permanentes (PMSM), aplicado al 
control de una articulación robótica de un grado de libertad.
El proyecto busca integrar los conocimientos 
adquiridos en el espacio curricular de Automática y Máquinas 
Eléctricas mediante el modelado, simulación, diseño
y evaluación del rendimiento 
de un sistema de control automático para regular el 
movimiento y la posición de un eje mecánico.

Durante el desarrollo del trabajo, se construirá el modelo 
matemático completo del sistema físico, incluyendo 
los subsistemas mecánico, electromagnético y térmico, y se 
realizarán simulaciones utilizando Matlab/Simulink. 
A partir de estos modelos, se diseñará una estrategia de 
control vectorial con desacoplamiento de variables y 
un controlador externo de movimiento tipo PID, evaluando 
además la estabilidad, controlabilidad y 
observabilidad del sistema.

Los resultados obtenidos a partir de las simulaciones 
permitirán verificar la capacidad del controlador para 
regular eficazmente la velocidad y el torque del motor, 
así como su robustez frente a perturbaciones 
externas y variaciones en los parámetros del sistema. 
También se analizarán los efectos térmicos en el 
motor y el impacto del torque de carga en distintas 
condiciones de operación.

Posteriormente, se reemplazarán los modelos ideales de 
sensores e inversor por modelos más realistas, 
incluyendo limitaciones como la saturación y un ancho de 
banda finito. Se analizará cómo estas restricciones 
afectan el desempeño del sistema de control.

Finalmente, se desarrollará una versión del controlador 
en tiempo discreto, evaluando la influencia de 
la frecuencia de muestreo sobre la precisión del sistema 
y su estabilidad, acercando la simulación a 
un escenario realista.


\newpage
\section{Introducción}
Este trabajo tiene como objetivo desarrollar un sistema de control automático para un accionamiento eléctrico de corriente alterna, basado en una máquina síncrona de imanes permanentes (PMSM), que actúa sobre la articulación de un brazo manipulador robótico con un solo grado de libertad. El sistema está diseñado para operar en un plano vertical, donde el eje de rotación se encuentra fijo a una base y sometido a la acción constante de la gravedad. La dinámica del sistema varía según las condiciones de carga útil aplicadas en el extremo del brazo.

El accionamiento propuesto está compuesto por un reductor de velocidad de engranajes planetarios, un motor eléctrico PMSM, y un conjunto de sensores encargados de proporcionar retroalimentación sobre la posición angular del motor, las corrientes de fase y la temperatura del estator. Con base en las especificaciones dadas, se desarrollará un modelo dinámico que integre los efectos mecánicos, electromagnéticos y térmicos, considerando tanto el comportamiento no lineal como su aproximación lineal, con el propósito de facilitar el análisis y diseño de estrategias de control.

El proyecto comprende distintas etapas, que incluyen el modelado del sistema físico, la simulación en lazo abierto, el diseño de un controlador vectorial mediante un modulador de torque equivalente, y la implementación de un controlador externo de movimiento en cascada. Estas simulaciones permitirán analizar la estabilidad, controlabilidad y observabilidad del sistema bajo distintas condiciones de operación. Además, se evaluará el desempeño térmico del motor y el impacto que tienen las limitaciones físicas de sensores e inversores sobre el comportamiento general del sistema.

% \begin{figure}[ht]
%     \centering
%     \includegraphics[width=1\textwidth]{imagenes/sistema_fisico.png}
%     \caption{Robot manipulador elemental de 1 g.d.l. en plano vertical (péndulo rígido actuado)}
% \end{figure}

\newpage
\section{Modelado, Análisis y Simulación dinámica del sistema físico a lazo abierto}

\subsection{Subsistema mecánico}

\[
-\left(\frac{2J_{hd+hEb}+2J_{hm+hb}\cdot i_h^{\,2}}{r_{hd}}\right)\ddot{l}_h(t)
=
\left(\frac{2b_{hd}+2b_{hm}\cdot i_h^{\,2}}{r_{hd}}\right)\cdot l_h(t)
+i_h\left(T_{hm}+T_{hb}\right)+T_{hEb}-F_{hw}r_{hd}
\]

\[
J_{eq}=\frac{J_{hd+hEb}}{i_h}+J_{hm+hb}\cdot{i}_h
\]

\[
B_{eq}=\frac{b_{hd}}{i_h}+b_{hm}\cdot i_h
\]

\end{document}

